\documentclass[conference]{IEEEtran}
\IEEEoverridecommandlockouts

\usepackage{amsmath,amssymb,amsfonts}
\usepackage{graphicx}
\usepackage{textcomp}
\usepackage{xcolor}
\usepackage[hidelinks]{hyperref}
\usepackage{booktabs}
\usepackage{array}

\begin{document}

\title{Multi-Agent Task Allocation with Deep Learning}

\author{\IEEEauthorblockN{Daria Shushkova}
\IEEEauthorblockA{Deep Learning -- Winter Term 2026/2027\\
Hamm-Lippstadt University of Applied Sciences\\
Lippstadt, Germany\\
daria.shushkova@stud.hshl.de}}

\maketitle

\begin{abstract}
Multi-Agent Task Allocation (MATA) is an important problem in multi-agent systems, where several autonomous agents must
coordinate their actions to complete tasks efficiently. Task allocation becomes more difficult in dynamic environments,
where new tasks may appear, agents have limited information, and some tasks require cooperation. Traditional approaches
use optimization algorithms, auction-based methods, and predefined rules to assign tasks. Recent advances in Deep
Reinforcement Learning (DRL) allow agents to learn task allocation strategies through interaction with their environment.

This paper examines the application of Deep Learning to dynamic multi-agent task allocation, focusing on decentralized
coordination and coalition formation. The theoretical foundations of Multi-Agent Reinforcement Learning (MARL) are
discussed, followed by a case study of MARL-MRT$^2$A, which combines Multi-Agent Proximal Policy Optimization (MAPPO) with
a U-Net architecture. The approach is examined using the GyMRT$^2$A simulation environment. An experimental reproduction
investigates how task allocation performance changes during training by evaluating policies at different training stages.
The evaluation considers episodic reward and task completion, with the market-based Project Manager-oriented Coalition
Formation Algorithm (PCFA) serving as a comparison baseline.
\end{abstract}

\section{Introduction}
Multi-agent systems consist of several autonomous agents that work in a shared environment to achieve individual or
common goals. Each agent collects information about its surroundings through sensors and observations. Agents can
exchange this information to coordinate their actions, avoid unnecessary work, and use available resources more
efficiently. However, effective cooperation also requires deciding which agent should perform each task. This problem
is known as Multi-Agent Task Allocation (MATA).

Task planning in multi-agent systems includes three main problems: task decomposition, task allocation, and task
scheduling \cite{tao2026survey}. Task decomposition divides complex tasks into smaller ones, task allocation assigns
tasks to different agents, and task scheduling determines the order in which tasks should be completed. Task allocation
is particularly important because agents may have different capabilities, locations, and available resources \cite{tao2026survey}.
In multi-robot systems, some tasks may also require several robots to work together. Bezerra et al. illustrate this
problem using emergency response, where multiple firetrucks must cooperate to handle incidents that cannot be managed
by one vehicle alone \cite{bezerra2025coalition}.

Task allocation becomes more difficult in dynamic environments, where new tasks may appear at any time and agents have
limited information about their surroundings. For example, an emergency incident may require multiple rescue robots,
while a smaller task can be completed by a single robot. Traditional task allocation approaches include optimization
algorithms, auction-based methods, and predefined rules \cite{tao2026survey}. More recent approaches use Deep
Reinforcement Learning (DRL), which allows agents to learn task allocation strategies through interaction with their
environment \cite{dai2025heterogeneous,bezerra2025coalition}.

Communication is an important part of decentralized task allocation. By sharing information about their observations,
selected tasks, and planned actions, agents can make better decisions and coordinate their work. However, communication
may become unreliable due to limited network coverage, interference, or message loss. Missing information can lead to
repeated task assignments, delays, and problems when several agents need to complete a task together. These
challenges highlight the importance of adaptive coordination strategies in dynamic multi-agent environments \cite{verma2025communication,yang2025priors}.

The practical part of this paper focuses on how Deep Reinforcement Learning can support cooperation between agents in
dynamic task allocation. MARL-MRT$^2$A is used as a case study to examine how task allocation performance develops during
training, particularly for tasks requiring different levels of cooperation \cite{bezerra2025coalition}. The reproduced
results are compared across training stages and considered alongside the traditional PCFA approach.

\section{Problem Formulation}
\begin{itemize}
    \item Agents and tasks in a dynamic environment.
    \item Task assignment and coalition formation.
    \item Challenges of limited information and changing tasks.
\end{itemize}

\section{Deep Learning-Based Task Allocation}
\begin{itemize}
    \item Deep Learning approaches to task allocation.
    \item Deep Reinforcement Learning and Multi-Agent Reinforcement Learning.
    \item Learning task selection and coordination strategies.
\end{itemize}

\section{Use Case and Experimental Setup}
\begin{itemize}
    \item Multi-robot task allocation in emergency response.
    \item MARL-MRT$^2$A and the GyMRT$^2$A simulation environment.
    \item Experimental evaluation and comparison with PCFA.
\end{itemize}

\section{Results and Evaluation}
\begin{itemize}
    \item How performance changes during training (comparison of saved model checkpoints).
    \item Task completion and reward results.
    \item Comparison with this Deep Learning approach with the classical PCFA.
\end{itemize}

\section{Discussion}
\begin{itemize}
    \item Advantages and limitations of Deep Learning for task allocation.
    \item Challenges of applying the approach in real environments.
\end{itemize}

\section{Conclusion}

\bibliographystyle{IEEEtran}
\bibliography{mata}

\end{document}
